% quark_scan_assumptions.tex
% Converted from quark_scan_assumptions.md — preserves every formula, table,
% assumption, numerical value, and footnote from the original document.
\documentclass[11pt]{article}

\usepackage[margin=1in]{geometry}
\usepackage{amsmath,amssymb}
\usepackage{booktabs}
\usepackage{xcolor}
\usepackage{hyperref}
\usepackage{enumitem}
\usepackage{array}
\usepackage{longtable}

\hypersetup{
  colorlinks=true,
  linkcolor=blue,
  citecolor=blue,
  urlcolor=blue
}

\newcommand{\bra}[1]{\langle #1 |}
\newcommand{\ket}[1]{| #1 \rangle}
\newcommand{\Mkk}{M_{\mathrm{KK}}}
\newcommand{\LamIR}{\Lambda_{\mathrm{IR}}}
\newcommand{\Mpl}{M_{\mathrm{Pl}}}
\newcommand{\vev}{v}
\newcommand{\gs}{g_s^*}

\title{Quark-Sector Scan: Complete Assumptions and Pipeline}
\author{}
\date{}

\begin{document}
\maketitle

This document traces every assumption, formula, and convention that goes into
producing the publication figures \texttt{fig1\_exclusion\_boundaries} and
\texttt{fig2\_mkk\_bound\_2007\_vs\_modern}. It follows the computation in the
order it actually executes, from geometry definition through to the final pixel.

%% ====================================================================
\section{Warped Geometry}
\label{sec:step0}

The spacetime is a slice of $\mathrm{AdS}_5$ (Randall--Sundrum~I) compactified
on an $S^1/\mathbb{Z}_2$ orbifold, bounded by two 3-branes: the UV brane at
the orbifold fixed point $y = 0$ (Planck-scale physics), and the IR brane at
$y = \pi r_c$ (TeV-scale physics).

The 5D metric in conformal coordinates (mostly-plus signature) is:
\begin{equation}
  ds^2 = (kz)^{-2}\bigl(\eta_{\mu\nu}\,dx^\mu\,dx^\nu + dz^2\bigr)
\end{equation}
where $\eta = \mathrm{diag}(-1,+1,+1,+1)$ and $z$ runs from $z_h = 1/k$
(UV brane) to $z_v = 1/\LamIR$ (IR brane).  In proper-distance coordinates:
\begin{equation}
  ds^2 = e^{-2k|y|}\,\eta_{\mu\nu}\,dx^\mu\,dx^\nu + dy^2
\end{equation}

Two parameters fix the entire geometry:

\begin{table}[h]
\centering
\begin{tabular}{@{}llll@{}}
\toprule
Parameter & Symbol & Default & Meaning \\
\midrule
AdS curvature & $k$ & $1.2209\times10^{19}$~GeV $(= \Mpl)$ & Sets the UV scale \\
IR scale & $\LamIR$ & scanned (500--20\,000~GeV) & Sets the TeV-scale brane \\
\bottomrule
\end{tabular}
\end{table}

Everything else follows from these two. The warp factor is
$\varepsilon = \LamIR / k$ (typically $\sim 10^{-15}$ for
$\LamIR \sim$ few~TeV). The orbifold radius is
$r_c = \ln(k/\LamIR)/(\pi k)$, related to the warp factor by
$\varepsilon = e^{-\pi k r_c}$. In conformal coordinates the UV brane sits at
$z_h = 1/k \sim 10^{-19}\;\mathrm{GeV}^{-1}$ and the IR brane at
$z_v = 1/\LamIR \sim 10^{-4}\;\mathrm{GeV}^{-1}$.

The electroweak VEV is $v = 174$~GeV throughout. This is the Higgs VEV, not
the reduced VEV $v/\sqrt{2} = 246/\sqrt{2}$; the code uses 174~GeV
consistently.

We take the Higgs to be localized on (or very near) the IR brane, so that the
Yukawa interaction is proportional to the fermion zero-mode value at the IR
brane and TeV-scale new physics is connected to the IR scale. We fix
$k = \Mpl = 1.2209\times10^{19}$~GeV and do not vary it; the warp factor
$\varepsilon$ and the hierarchy are controlled entirely by $\LamIR$.
Backreaction of bulk fields on the geometry is neglected, and the metric is
taken to be the pure $\mathrm{AdS}_5$ slice throughout.

%% ====================================================================
\section{Zero-Mode Profiles ($f$-factors)}
\label{sec:step1}

Each 5D fermion field has a bulk Dirac mass $M_5$.  The dimensionless bulk mass
parameter is $c = M_5/k$.  The zero-mode wavefunction in the extra dimension is
determined by~$c$: it satisfies the 5D Dirac equation with the AdS background,
and the boundary conditions select a single normalizable zero mode per
chirality.

The physical quantity that enters all 4D couplings is the value of this
normalized zero-mode wavefunction at the IR brane.  For left-handed zero modes
(the ones relevant for Higgs overlap), this ``$f$-factor'' is:
\begin{align}
  f_{\mathrm{IR}}^2(c,\varepsilon) &= \frac{\tfrac{1}{2}-c}
    {1 - \varepsilon^{1-2c}}\,, \\
  f_{\mathrm{IR}}(c,\varepsilon) &= \sqrt{\max\!\bigl(f_{\mathrm{IR}}^2,\;0\bigr)}\,.
\end{align}

The formula arises from integrating the square of the zero-mode profile
$z^{2-c}$ over the extra dimension to normalize it, then evaluating at
$z = z_v$.  The $\max(\ldots,0)$ is a numerical safeguard against round-off.

The three regimes are:
\begin{itemize}
  \item $c > 1/2$: The exponent $(1-2c)$ is negative, so
    $\varepsilon^{1-2c}$ is huge (recall $\varepsilon\sim10^{-15}$).  The
    denominator is dominated by $-\varepsilon^{1-2c}$, and
    $f_{\mathrm{IR}}^2 \sim (c-\tfrac{1}{2})\,\varepsilon^{2c-1}$, which is
    exponentially small.
    \emph{These fermions are UV-localized and have tiny overlap with the
    IR-brane Higgs.}

  \item $c < 1/2$: The exponent $(1-2c)$ is positive, so
    $\varepsilon^{1-2c}\sim 0$.  Then $f_{\mathrm{IR}}^2 \sim (1/2-c)$,
    which is $O(1)$.
    \emph{These fermions are IR-localized and have large Yukawa overlap.}

  \item $c = 1/2$: The formula is $0/0$.  The limit is
    $f_{\mathrm{IR}}^2 = 1/(-2\ln\varepsilon)$.
    For $\varepsilon\sim10^{-15}$, this gives
    $f_{\mathrm{IR}}^2 \sim 1/69 \sim 0.014$.
\end{itemize}

At $\varepsilon = 3000/1.2209\times10^{19} \approx 2.46\times10^{-16}$, for
instance, $c = 0.72$ (light quark) gives $f_{\mathrm{IR}} \sim 4.4\times10^{-8}$,
$c = 0.51$ (near flat) gives $f_{\mathrm{IR}} \sim 0.11$, and $c = 0.30$
(heavy quark) gives $f_{\mathrm{IR}} \sim 0.45$. That is 10 orders of magnitude
of dynamic range between $c = 0.72$ and $c = 0.30$; the mass hierarchies come
from geometry, not from hierarchical Yukawa couplings.

We use $c = M_5/k$, so positive $c$ means the bulk mass
and the AdS curvature have the same sign. This is the standard convention in the
RS flavor literature (Agashe, Contino, Perez 2005; Csaki et al.\ 2009).

There is also a UV-brane overlap $f_{\mathrm{UV}}$, used for neutrino Majorana
masses in the lepton sector, but it does not enter the quark scan.

%% ====================================================================
\section{Scan Grid}
\label{sec:step2}

The dense scan (the one used for the publication figures) sweeps a 3D grid.
Here \texttt{overall\_scale} is a common multiplicative prefactor on both $Y_u$
and $Y_d$ (see the Yukawa section below); it controls the overall size of the
Yukawa entries and therefore the tradeoff between quark masses and FCNCs.
Different values of \texttt{overall\_scale} lead to different Yukawa structures
even after the fitter re-optimizes to hit the same quark masses: if
\texttt{overall\_scale} is large, the fitter drives the singular values small to
compensate, and vice versa. These different internal configurations produce
different off-diagonal entries in the mass basis and therefore different FCNC
predictions. Sweeping \texttt{overall\_scale} maps out this family of
solutions. Over the scan range 1.5 to 6.0, the largest singular value
(top-like) runs from $\sim 1.5$ to $\sim 6$ and the smallest (up-like) from
$\sim 0.005$ to $\sim 0.02$. At the high end the top Yukawa is well above the
naive $O(1)$ expectation.

\begin{table}[h]
\centering
\begin{tabular}{@{}llll@{}}
\toprule
Parameter & Grid & Range & Points \\
\midrule
$r$ & log-uniform & 0.02--2.0 & 100 \\
\texttt{overall\_scale} & linear & 1.5--6.0 & 20 \\
$\LamIR$ & log-uniform & 500--20\,000~GeV & 50 \\
\bottomrule
\end{tabular}
\end{table}

Total: $100 \times 20 \times 50 = 100{,}000$ scan points.

Each point is fitted independently from the same deterministic seed (no
warm-starting between points).  Points are fit in parallel; the result is
independent of execution order.

At scan time, several parameters are held fixed. The KK-mass convention is
$\xi_{\mathrm{KK}} = 1.0$, meaning $\Mkk = \LamIR$ in the scan's bookkeeping;
the Wilson coefficients use $\Mkk = \LamIR$ as the propagator mass. (See
\hyperref[sec:step16b]{Step~16b} for the post-hoc conversion to a physical
KK-gluon mass.)

\textcolor{red}{\textbf{To do:} Future scans should set
$\xi_{\mathrm{KK}} = 2.4487$ (the physical first KK-gluon mass from the
gauge-sector Bessel root) so that $\Mkk = m_g^{(1)}$ from the start.  The
current post-hoc correction (\hyperref[sec:step16b]{Step~16b}) rescales
the $1/\Mkk^2$ propagator factor but does not account for the fact that
the QCD running of Wilson coefficients (Step~13) also starts at the wrong
matching scale.  Using the physical $\xi_{\mathrm{KK}}$ would eliminate
both issues.}

The enhanced KK-gluon coupling is $\gs = 3.0$ (not the perturbative $g_s$),
taken from the one-loop estimate of Gedalia, Grossman, Nir, Perez (2009). It
enters the coupling computation at scan time, so all stored Wilson coefficients
and ratios already include this enhancement.

\textcolor{red}{\textbf{To do:} Confirm the correct value of $\gs$ with
Lisa.  The tree-level estimate is $\gs \sim 6$ (Csaki, Falkowski, Weiler
2008); the one-loop value is $\gs \sim 3$ (Gedalia et al.\ 2009).  Since
$\gs$ enters as ${\gs}^2$, the difference between these two choices is a
factor of~4 in all Wilson coefficients and constraint ratios.}

The remaining fixed inputs are $k = 1.2209\times10^{19}$~GeV $(= \Mpl)$ and
$v = 174$~GeV.

The initial fit seed uses $O(1)$ singular values and small CKM-like
left-rotation angles, chosen by hand to be in the right ballpark for SM-like
quark masses.  Right rotations in the seed are set to identity during
canonicalization before the fit begins
(\hyperref[sec:step10]{Step~10}).

%% ====================================================================
\section{Yukawa Spurions}
\label{sec:step3}

$Y_u$ and $Y_d$ are complex $3\times3$ matrices.  They are the 5D Yukawa
coupling matrices that appear in the action on the IR brane as
$\mathcal{L} = Y_{ij}\,Q_i\,H\,u_j$ (schematic).
In the scan, they are parametrized via a singular value decomposition:
\begin{align}
  Y_u &= \texttt{overall\_scale} \times U_{\mathrm{left}}^u
    \cdot \operatorname{diag}(\sigma_u)
    \cdot \bigl(U_{\mathrm{right}}^u\bigr)^\dagger\,, \\
  Y_d &= \texttt{overall\_scale} \times U_{\mathrm{left}}^d
    \cdot \operatorname{diag}(\sigma_d)
    \cdot \bigl(U_{\mathrm{right}}^d\bigr)^\dagger\,,
\end{align}
where $\sigma_u = (\sigma_{u1},\,\sigma_{u2},\,\sigma_{u3})$ are 3 positive
singular values (ordered light-to-heavy) controlling the hierarchy within each
Yukawa, and $\sigma_d$ likewise for the down sector.
$U_{\mathrm{left}}^{u,d}$ are CKM-like $3\times3$ unitary matrices, each
parametrized by three Euler angles
$(\theta_{12},\,\theta_{13},\,\theta_{23})$ and one CP-violating phase
$(\delta)$. These control the orientation of the Yukawa in flavor space.
$U_{\mathrm{right}}^{u,d}$ are fixed to the $3\times3$ identity matrix
(not fitted), which restricts the parameter space (see below).
Finally, \texttt{overall\_scale} is a common real positive prefactor multiplying
both $Y_u$ and $Y_d$.

The unitary matrices use the standard CKM (PDG) parametrization:
\begin{equation}
  U = \begin{pmatrix}
    c_{12}\,c_{13}
      & s_{12}\,c_{13}
      & s_{13}\,e^{-i\delta} \\
    -s_{12}\,c_{23} - c_{12}\,s_{23}\,s_{13}\,e^{i\delta}
      & c_{12}\,c_{23} - s_{12}\,s_{23}\,s_{13}\,e^{i\delta}
      & s_{23}\,c_{13} \\
    s_{12}\,s_{23} - c_{12}\,c_{23}\,s_{13}\,e^{i\delta}
      & -c_{12}\,s_{23} - s_{12}\,c_{23}\,s_{13}\,e^{i\delta}
      & c_{23}\,c_{13}
  \end{pmatrix}
\end{equation}
where $c_{ij} = \cos\theta_{ij}$, $s_{ij} = \sin\theta_{ij}$.  All angles are
wrapped into $[-\pi,\pi)$.

Before the fit begins, \texttt{overall\_scale} is
absorbed into the singular values and set to~$1.0$; right rotations are set to
identity.  After canonicalization:
\begin{align}
  Y_u &= U_{\mathrm{left}}^u \cdot
    \operatorname{diag}(\sigma_u)\,, \qquad
  Y_d = U_{\mathrm{left}}^d \cdot
    \operatorname{diag}(\sigma_d)\,.
\end{align}

With right rotations set
to identity:
\begin{align}
  Y_u^\dagger Y_u &= \operatorname{diag}(\sigma_u)^2 \quad\text{(diagonal)}\,, \\
  Y_d^\dagger Y_d &= \operatorname{diag}(\sigma_d)^2 \quad\text{(diagonal)}\,, \\
  Y_u Y_u^\dagger &= U_{\mathrm{left}}^u \cdot \operatorname{diag}(\sigma_u)^2
    \cdot \bigl(U_{\mathrm{left}}^u\bigr)^\dagger
    \quad\text{(generally dense)}\,, \\
  Y_d Y_d^\dagger &= U_{\mathrm{left}}^d \cdot \operatorname{diag}(\sigma_d)^2
    \cdot \bigl(U_{\mathrm{left}}^d\bigr)^\dagger
    \quad\text{(generally dense)}\,.
\end{align}

So $C_u$ and $C_d$ (defined next) are automatically diagonal, while $C_Q$ is
generally dense.  The dominant source of flavor misalignment in this scan is
the left-handed matrix $C_Q$, built from the combination of the two left
rotations $U_{\mathrm{left}}^u$ and $U_{\mathrm{left}}^d$. (Right-handed
off-diagonal couplings can still appear after the SVD rotation to the mass
basis, but the restriction to diagonal $C_u$, $C_d$ removes much of the
freedom a general scan would have.)

We freeze right rotations to identity, which restricts the Yukawa space to 15
free parameters instead of 37 for two general complex $3\times 3$ matrices.

%% ====================================================================
\section{Bulk Mass Matrices}
\label{sec:step4}

From the Yukawa spurions, three Hermitian matrices are constructed:
\begin{align}
  C_u &= Y_u^\dagger \cdot Y_u
    &&\text{[controls RH up-type quark localization]}\,, \\
  C_d &= Y_d^\dagger \cdot Y_d
    &&\text{[controls RH down-type quark localization]}\,, \\
  C_Q &= r \cdot (Y_u\,Y_u^\dagger) + (Y_d\,Y_d^\dagger)
    &&\text{[controls LH doublet localization]}\,.
\end{align}

In the RS NMFV framework, the 5D
bulk mass matrix for the left-handed quark doublets is assumed to be a function
of the Yukawa spurions (since the Yukawas are the only sources of flavor
breaking).  At leading order in the spurion expansion:
\begin{equation}
  C_Q = \alpha_Q\,\mathbf{I} + \beta_Q\,Y_d Y_d^\dagger
    + \gamma_Q\,Y_u Y_u^\dagger + \cdots
\end{equation}

The parameter $r = \gamma_Q/\beta_Q$ measures the relative weight of the
up-type vs down-type Yukawa contribution.  The code uses the simplified form
with $\beta_Q = 1$ (normalized out) and no flavor-universal $\alpha_Q$ term:
\begin{equation}
  C_Q = r\cdot(Y_u\,Y_u^\dagger) + (Y_d\,Y_d^\dagger)
\end{equation}

$C_Q$ is Hermitian and positive semi-definite (it is a sum of positive
semi-definite matrices with non-negative coefficients).

When $r = 0$, $C_Q = Y_d\,Y_d^\dagger$.  Working
in the basis where $Y_d$ is diagonal (which we can always choose), $C_Q$ is
also diagonal.  In that basis the left-handed overlap matrix $F_Q$ is diagonal,
and there are no tree-level FCNCs in the down sector.  This is the ``alignment''
limit.

When $r > 0$, the $Y_u\,Y_u^\dagger$ term introduces off-diagonal entries
(since $Y_u$ contains CKM-like rotations relative to $Y_d$).  These
off-diagonal entries in $C_Q$ misalign the left-handed quark basis from the
down-Yukawa eigenbasis, generating tree-level FCNCs.

The literature
(Csaki et al.\ 0907.0474) writes
$C_Q = \alpha_Q\,\mathbf{I} + \beta_Q\,Y_d Y_d^\dagger +
\gamma_Q\,Y_u Y_u^\dagger$.
The $\alpha_Q\,\mathbf{I}$ term provides a flavor-universal baseline bulk mass
(typically $\alpha_Q \sim 0.5$--$0.7$ to localize quarks near the UV brane).
Our scan omits $\alpha_Q\,\mathbf{I}$ entirely.  In its place, an ad hoc
nonlinear map (\hyperref[sec:step5]{Step~5}) clamps the eigenvalues of $C_Q$
into a bounded window $[0.30,\,0.72]$ after diagonalization. This is not
equivalent to including $\alpha_Q$: $\alpha_Q\,\mathbf{I}$ is a linear
flavor-universal matrix term in the spurion expansion, whereas the sigmoid is
a nonlinear spectral remap that preserves eigenvectors but distorts eigenvalue
splittings in a way that has no counterpart in the MFV expansion.

\textcolor{red}{\textbf{To do:} The note from Jackie prescribes the explicit
$(\alpha_Q,\,\beta_Q,\,\gamma_Q)$ parametrization with $C_Q =
\alpha_Q\,\mathbf{I} + \beta_Q\,Y_d Y_d^\dagger + \gamma_Q\,Y_u Y_u^\dagger$,
where all three coefficients are free $O(1)$ parameters.  The current scan uses
a different structure (no $\alpha_Q$, nonlinear BulkMassMap).  A future scan
following the note should implement the explicit parametrization directly.}

Because $U_{\mathrm{right}} = I$ (\hyperref[sec:step3]{Step~3}), both $C_u$ and
$C_d$ are diagonal: $C_u = \operatorname{diag}(\sigma_u^2)$,
$C_d = \operatorname{diag}(\sigma_d^2)$.  Only $C_Q$ carries off-diagonal
structure.

%% ====================================================================
\section{Diagonalization and Bulk Mass Mapping}
\label{sec:step5}

\subsection{5a: Hermitian eigendecomposition}
\label{sec:step5a}

Each of $C_Q$, $C_u$, $C_d$ is diagonalized (Hermitian eigendecomposition):
\begin{equation}
  C_X = \text{rotation}_X \cdot \operatorname{diag}(\lambda_X)
    \cdot \text{rotation}_X^\dagger
\end{equation}
with eigenvalues sorted ascending (lightest generation first).  Since $C_u$ and
$C_d$ are already diagonal, their rotations are trivially the identity.  For
$C_Q$ the rotation is nontrivial and defines the basis transformation from the
spurion basis to the bulk-mass eigenbasis.

\subsection{5b: Nonlinear map from eigenvalues to bulk masses}
\label{sec:step5b}

The eigenvalues $\lambda_X$ (non-negative reals) are mapped to physical bulk
mass parameters $c_X$ through a bounded monotone sigmoid:
\begin{equation}
  c_X = c_{\mathrm{uv}} - (c_{\mathrm{uv}} - c_{\mathrm{ir}})
    \cdot \frac{\lambda}{\lambda + \text{eigen\_scale}}
\end{equation}

With the defaults $c_{\mathrm{uv}} = 0.72$, $c_{\mathrm{ir}} = 0.30$,
$\text{eigen\_scale} = 1.0$, this becomes:
\begin{equation}
  c_X = 0.72 - 0.42\;\frac{\lambda}{\lambda + 1}
\end{equation}

This is a monotone map bounded in $[0.30,\,0.72]$:
$\lambda = 0 \Rightarrow c = 0.72$ (UV-localized),
$\lambda = 1 \Rightarrow c = 0.51$,
$\lambda \to \infty \Rightarrow c \to 0.30$ (IR-localized).
The map enforces that all bulk masses stay in a physically sensible range
($c < 0$ gives imaginary $f$-factors; $c > 0.8$ gives essentially zero overlap).

The window $[0.30,\,0.72]$ is a design choice.
$c_{\mathrm{uv}} = 0.72$ means that a fermion with zero spurion eigenvalue
(i.e., decoupled from the Yukawa structure) sits at $c = 0.72$, which is
UV-localized and very light ($f_{\mathrm{IR}} \sim 4\times10^{-8}$).
$c_{\mathrm{ir}} = 0.30$ means that the most IR-localized fermion has
$c = 0.30$ ($f_{\mathrm{IR}} \sim 0.45$).  The top quark, which needs
$f_{\mathrm{IR}} \sim O(1)$, would need $c \sim 0.3$ or below; the BulkMassMap
cannot produce $c < 0.30$.  This forces the top Yukawa singular value to be
$\sim 2.4$ to reproduce $m_t = 172$~GeV (see \hyperref[sec:step8]{Step~8}),
which is perturbative but larger than the ``anarchic $O(1)$'' expectation.

The same sigmoid is applied to all three sectors ($Q$, $u$, $d$) with the
same parameters.  This removes sector-dependent universal offsets (i.e.\
independent $\alpha_Q$, $\alpha_u$, $\alpha_d$) and sign/cancellation freedom
in the $O(1)$ coefficients.  It is not equivalent to the standard NMFV
parametrization: $\alpha_Q\,\mathbf{I}$ is a linear flavor-universal matrix
term, whereas the sigmoid is a post-diagonalization spectral remap that
preserves eigenvectors but distorts eigenvalue splittings nonlinearly.  In the
standard parametrization, $\alpha_Q$ provides the baseline and the eigenvalues
of $C_Q$ are directly the bulk masses, with no remapping needed.

\textcolor{red}{\textbf{To do:} The ad hoc sigmoid is the assumption most
likely to bias both the top-sector masses and the FCNC bounds.  With the
explicit $(\alpha_Q,\,\beta_Q,\,\gamma_Q)$ parametrization (and analogous
$\alpha_u$, $\alpha_d$ for the RH sectors), it would be replaced by a direct
identification $c_{X,i} = \text{eigenvalue}_i(C_X)$, with independent baselines
per sector.}

%% ====================================================================
\section{Zero-Mode Overlaps}
\label{sec:step6}

With the physical bulk masses $c_Q$, $c_u$, $c_d$ (each a length-3 vector, one
per generation), the zero-mode IR-brane overlaps are computed element-wise using
the formula from \hyperref[sec:step1]{Step~1}:
\begin{align}
  F_Q[i] &= f_{\mathrm{IR}}(c_Q[i],\,\varepsilon)
    \qquad\text{for } i = 1,2,3 \text{ (generations)}\,, \\
  F_u[i] &= f_{\mathrm{IR}}(c_u[i],\,\varepsilon)\,, \\
  F_d[i] &= f_{\mathrm{IR}}(c_d[i],\,\varepsilon)\,.
\end{align}

These are real, positive numbers stored as 1D arrays of length~3.  They will be
used as diagonal matrices in the mass formula.

Because of the BulkMassMap range $[0.30,\,0.72]$, all $c$ values are $> 1/2$
except possibly the third generation.  In practice:
\begin{itemize}
  \item First-generation quarks ($c \sim 0.7$):
    $F \sim 10^{-7}$~to~$10^{-8}$
  \item Second-generation quarks ($c \sim 0.55$--$0.65$):
    $F \sim 10^{-3}$~to~$10^{-5}$
  \item Third-generation quarks ($c \sim 0.30$--$0.45$):
    $F \sim 0.1$~to~$0.45$
\end{itemize}

%% ====================================================================
\section{Rotation to the Bulk-Mass Basis}
\label{sec:step7}

The diagonalization of $C_Q$, $C_u$, $C_d$ in \hyperref[sec:step5]{Step~5}
defines three new bases.  The Yukawas must be rotated into the bulk-mass
eigenbasis before the mass matrices can be built, because the overlap matrices
$F_Q$, $F_u$, $F_d$ are diagonal only in that basis.

The rotation is:
\begin{align}
  Y_u^{\mathrm{bulk}} &= \text{rotation}_Q^\dagger \cdot Y_u
    \cdot \text{rotation}_u\,, \\
  Y_d^{\mathrm{bulk}} &= \text{rotation}_Q^\dagger \cdot Y_d
    \cdot \text{rotation}_d\,,
\end{align}
where $\text{rotation}_Q$, $\text{rotation}_u$, $\text{rotation}_d$ are the
unitary eigenvector matrices from \hyperref[sec:step5]{Step~5}.

This rotation does \emph{not} diagonalize the Yukawas; it diagonalizes the
bulk mass matrices.  The rotated Yukawas $Y_u^{\mathrm{bulk}}$ and
$Y_d^{\mathrm{bulk}}$ are generally dense $3\times3$ complex matrices.
Off-diagonal entries in these rotated Yukawas, combined with the non-degenerate
diagonal overlap factors, produce off-diagonal entries in the mass matrices and
drive CKM mixing.

(Since $C_u$ and $C_d$ are already diagonal, $\text{rotation}_u$ and
$\text{rotation}_d$ reduce to the identity in practice.)

%% ====================================================================
\section{4D Mass Matrices}
\label{sec:step8}

The effective 4D quark mass matrices in the bulk-mass eigenbasis are:
\begin{align}
  M_u &= 2v \cdot \operatorname{diag}(F_Q) \cdot Y_u^{\mathrm{bulk}}
    \cdot \operatorname{diag}(F_u)\,, \\
  M_d &= 2v \cdot \operatorname{diag}(F_Q) \cdot Y_d^{\mathrm{bulk}}
    \cdot \operatorname{diag}(F_d)\,,
\end{align}
where $v = 174$~GeV is the electroweak VEV.

With $v = 174$~GeV, the prefactor
$2v = 348$~GeV.  Note that $2\times174 = 348 \sim \sqrt{2}\times246$, where
246~GeV is the standard electroweak VEV
($v_{\mathrm{SM}} = 246$~GeV, with the Higgs field normalization
$\langle H \rangle = v_{\mathrm{SM}}/\sqrt{2} = 174$~GeV).  So the mass
formula is equivalently:
\begin{equation}
  M = \sqrt{2}\;v_{\mathrm{SM}} \cdot \operatorname{diag}(F_Q)
    \cdot Y^{\mathrm{bulk}} \cdot \operatorname{diag}(F)
\end{equation}

The factor of $\sqrt{2}$ arises from the normalization of the Higgs doublet VEV
relative to the mass term in the Lagrangian.  The 5D Yukawa couplings $Y$ in
this code are dimensionless (they absorb the factor of $k$ that appears in the
lepton-sector formula $m_E = 2vk\,f_L\,Y_E\,f_E$ used elsewhere in this
repo).  What matters is that the target quark masses
(\hyperref[sec:step10]{Step~10}) are defined consistently with this convention.

The left overlap
$\operatorname{diag}(F_Q)$ is the same for both sectors (since $Q$ is the
$SU(2)$ doublet containing both $u_L$ and $d_L$).  The right overlaps
$\operatorname{diag}(F_u)$ and $\operatorname{diag}(F_d)$ differ.  The Yukawas
$Y_u^{\mathrm{bulk}}$ and $Y_d^{\mathrm{bulk}}$ are generally dense.

The physical quark mass for generation $i$ is roughly:
\begin{equation}
  m_i \sim 2v \cdot F_Q[i] \cdot Y^{\mathrm{bulk}}_{ii}
    \cdot F_{u\,\text{or}\,d}[i]
\end{equation}
for the approximately diagonal case.  The exponential hierarchy in the
$F$-factors produces the quark mass hierarchy without requiring hierarchical
Yukawa entries.

%% ====================================================================
\section{Physical Masses and CKM}
\label{sec:step9}

Each mass matrix is decomposed via SVD (biunitary diagonalization) using the
convention:
\begin{equation}
  M = U_L \cdot \operatorname{diag}(\text{singular values})
    \cdot U_R^\dagger
\end{equation}
so that $U_L^\dagger \cdot M \cdot U_R =
\operatorname{diag}(m_1,\,m_2,\,m_3)$ with $m_i \ge 0$, sorted ascending
($m_1 < m_2 < m_3$ corresponds to $u < c < t$ or $d < s < b$).

The CKM matrix is constructed from the left-handed rotation matrices:
\begin{equation}
  V_{\mathrm{CKM}} = (U_L^u)^\dagger \cdot U_L^d
\end{equation}

This is the standard definition: $V_{\mathrm{CKM}}$ rotates from the down mass
eigenbasis to the up mass eigenbasis for left-handed quarks.

Four CKM observables are extracted:
\begin{enumerate}
  \item $|V_{us}| = |V_{\mathrm{CKM}}[0,1]|$ (Cabibbo angle)
  \item $|V_{cb}| = |V_{\mathrm{CKM}}[1,2]|$
  \item $|V_{ub}| = |V_{\mathrm{CKM}}[0,2]|$
  \item $J = \mathrm{Im}\!\bigl(V_{\mathrm{CKM}}[0,1] \cdot
    V_{\mathrm{CKM}}[1,2] \cdot V_{\mathrm{CKM}}[0,2]^* \cdot
    V_{\mathrm{CKM}}[1,1]^*\bigr)$,
    the Jarlskog invariant (rephasing-invariant measure of CP violation)
\end{enumerate}

%% ====================================================================
\section{The Fit}
\label{sec:step10}

Steps~3 through~9 are wrapped inside a Trust Region Reflective (TRF) optimizer
(\texttt{scipy.optimize.least\_squares}).  The fit adjusts the Yukawa
parameters until the predicted quark masses and CKM match the observed values.

\subsection{10a: Free parameters (14 total)}
\label{sec:step10a}

The optimizer works in a ``canonical chart'' with 14 real coordinates:

\begin{enumerate}
  \item[1--3.] $\ln(\text{physical\_}\sigma_u[1])$,
    $\ln(\text{physical\_}\sigma_u[2])$,
    $\ln(\text{physical\_}\sigma_u[3])$
    (log of the up-sector singular values; the log ensures positivity)
  \item[4--6.] $\ln(\text{physical\_}\sigma_d[1])$,
    $\ln(\text{physical\_}\sigma_d[2])$,
    $\ln(\text{physical\_}\sigma_d[3])$
    (log of the down-sector singular values)
  \item[7--10.] $\theta_{12}^u$, $\theta_{13}^u$, $\theta_{23}^u$, $\delta^u$
    (four parameters of the up-sector left rotation)
  \item[11--14.] $\theta_{12}^d$, $\theta_{13}^d$, $\theta_{23}^d$, $\delta^d$
    (four parameters of the down-sector left rotation)
\end{enumerate}

All angles are in $[-\pi,\pi)$ after canonicalization.

During the fit, $r$, \texttt{overall\_scale} (absorbed into singular
values as~$1.0$), $\LamIR$, $k$, $v$, the BulkMassMap parameters, and the
right rotations $(= I)$ are all held fixed.

\subsection{10b: Target observables (10 total)}
\label{sec:step10b}

\begin{table}[h]
\centering
\begin{tabular}{@{}lll@{}}
\toprule
Observable & Target value & Source \\
\midrule
$m_u$ & 0.0013~GeV & Repo-owned target set \\
$m_c$ & 0.62~GeV & Repo-owned target set \\
$m_t$ & 172.0~GeV & Repo-owned target set \\
$m_d$ & 0.0028~GeV & Repo-owned target set \\
$m_s$ & 0.057~GeV & Repo-owned target set \\
$m_b$ & 2.86~GeV & Repo-owned target set \\
$|V_{us}|$ & 0.225 & CKM parametrization \\
$|V_{cb}|$ & 0.0415 & CKM parametrization \\
$|V_{ub}|$ & 0.00361 & CKM parametrization \\
$J$ & $\sim 3.08\times10^{-5}$ &
  From $(\theta_{12}{=}0.2274,\;\theta_{13}{=}0.00368,$\\
  && $\theta_{23}{=}0.0415,\;\delta{=}1.196)$ \\
\bottomrule
\end{tabular}
\end{table}

The target CKM is built from a CKM-like unitary with angles
$(\theta_{12}=0.2274,\;\theta_{13}=0.00368,\;\theta_{23}=0.0415,\;
\delta=1.196)$.  The four CKM observables $(|V_{us}|,\;|V_{cb}|,\;|V_{ub}|,\;
J)$ are extracted from this matrix.

These target masses are repo-owned ``effective'' quark
masses, not $\overline{\mathrm{MS}}$ masses at any single consistent
renormalization scale.  Some values ($m_u = 1.3$~MeV, $m_d = 2.8$~MeV) are
close to $\overline{\mathrm{MS}}$ at $\mu = 2$~GeV; others
($m_c = 0.62$~GeV) are closer to $\overline{\mathrm{MS}}$ at
$\mu \sim 3$~TeV; $m_b = 2.86$~GeV is close to $\overline{\mathrm{MS}}$ at
$\mu = m_b$.  The values were chosen for scan stability and to give physically
reasonable bulk mass parameters, not for precise consistency with a single
renormalization scheme.  Since the fit determines the bulk masses (and thus the
$f$-factors and FCNC couplings) from these targets, any systematic shift in
the targets propagates into the constraint ratios, introducing unquantified
systematic uncertainty.

\textcolor{red}{\textbf{To do:} Replace these with properly defined PDG
$\overline{\mathrm{MS}}$ quark masses at a consistent renormalization scale
(e.g., all at $\mu = 2$~GeV, or all at $\mu = \Mkk$).  The current
inconsistency between the scales at which the target masses are defined is
an uncontrolled systematic.}

\subsection{10c: Residual vector (10 components)}
\label{sec:step10c}

The optimizer minimizes a vector of 10 residuals:
\begin{equation}
\text{residuals} = \begin{pmatrix}
  \ln(m_u^{\mathrm{pred}}/0.0013) \\[2pt]
  \ln(m_c^{\mathrm{pred}}/0.62) \\[2pt]
  \ln(m_t^{\mathrm{pred}}/172.0) \\[2pt]
  \ln(m_d^{\mathrm{pred}}/0.0028) \\[2pt]
  \ln(m_s^{\mathrm{pred}}/0.057) \\[2pt]
  \ln(m_b^{\mathrm{pred}}/2.86) \\[2pt]
  (|V_{us}|^{\mathrm{pred}} - 0.225) / 0.225 \\[2pt]
  (|V_{cb}|^{\mathrm{pred}} - 0.0415) / 0.0415 \\[2pt]
  (|V_{ub}|^{\mathrm{pred}} - 0.00361) / 0.00361 \\[2pt]
  (J^{\mathrm{pred}} - J^{\mathrm{target}}) / |J^{\mathrm{target}}|
\end{pmatrix}
\end{equation}

Mass residuals are in log-space (dimensionless); CKM residuals are fractional
deviations normalized by the target value.

The fit score is $\sqrt{\mathrm{mean}(\text{residuals}^2)}$, i.e.\ the RMS of
all 10 residuals.  A score of~$0.1$ means $\sim 10\%$ average deviation across
all observables.

\subsection{10d: Optimizer settings}
\label{sec:step10d}

The method is Trust Region Reflective (TRF) via
\texttt{scipy.optimize.least\_squares}, with a maximum of 120 evaluations (each
evaluation runs Steps~3--9 once).

Points with $\text{fit\_score} > 0.1$ after the optimizer finishes are flagged
as poorly converged and discarded at plotting time
(\hyperref[sec:step16a]{Step~16a}).

%% ====================================================================
\section{Mass-Basis KK-Gluon Couplings}
\label{sec:step11}

After the fit converges, we have: (1)~the diagonal overlap profiles $F_Q$,
$F_u$, $F_d$, and (2)~the unitary rotation matrices $U_L^u$, $U_R^u$,
$U_L^d$, $U_R^d$ from the SVD of the mass matrices.  These are combined to
build the KK-gluon coupling matrices in the physical quark mass basis.

We model the KK-gluon coupling to zero-mode quarks as proportional to
$\operatorname{diag}(F^2)$ in the bulk-mass eigenbasis (an overlap ansatz that
captures localization dependence but is not the full 5D overlap integral in
exact RS notation).  Multiplied by $\gs$ and rotated into the mass basis:
\begin{equation}
  (\text{overlap})_{ij}^{\text{mass basis}} = \bigl(U^\dagger \cdot
    \operatorname{diag}(F^2) \cdot U\bigr)_{ij}
\end{equation}
where $U$ is the appropriate left or right rotation matrix from the SVD.

Four coupling matrices are built:
\begin{align}
  g_L^{\mathrm{down}}[i,j] &= \gs \cdot \bigl(U_L^d\bigr)^\dagger
    \cdot \operatorname{diag}(F_Q^2) \cdot U_L^d \Big|_{ij}\,, \\
  g_R^{\mathrm{down}}[i,j] &= \gs \cdot \bigl(U_R^d\bigr)^\dagger
    \cdot \operatorname{diag}(F_d^2) \cdot U_R^d \Big|_{ij}\,, \\
  g_L^{\mathrm{up}}[i,j] &= \gs \cdot \bigl(U_L^u\bigr)^\dagger
    \cdot \operatorname{diag}(F_Q^2) \cdot U_L^u \Big|_{ij}\,, \\
  g_R^{\mathrm{up}}[i,j] &= \gs \cdot \bigl(U_R^u\bigr)^\dagger
    \cdot \operatorname{diag}(F_u^2) \cdot U_R^u \Big|_{ij}\,.
\end{align}

Each is a Hermitian $3\times3$ matrix.  The diagonal entries are the
flavor-conserving couplings; the off-diagonal entries drive FCNCs.

Note that $F_Q$ appears in both the LH-down and LH-up couplings because both
$u_L$ and $d_L$ live in the same $SU(2)$ doublet $Q$, which has the same bulk
mass and therefore the same overlap profile.  The difference between
$g_L^{\mathrm{down}}$ and $g_L^{\mathrm{up}}$ comes entirely from the
different rotation matrices $U_L^d$ vs $U_L^u$.

The code uses $\gs = 3.0$ (passed at scan time),
replacing the perturbative QCD coupling
$g_s = \sqrt{4\pi\alpha_s(\Mkk)}$.  In RS, the first KK-gluon has a
wavefunction peaked near the IR brane, enhancing its coupling to IR-localized
fermions by a factor $\sqrt{2\,k\,\pi\,r_c} \sim 6$--$8$ at tree level.  The
one-loop corrected value is $\gs \sim 3$ (Gedalia et al.\ 2009).

Since $\gs$ enters the Wilson coefficients as ${\gs}^2$, the difference between
$\gs = 1$ (perturbative) and $\gs = 3$ is a factor of~$9$ in all constraint
ratios.

%% ====================================================================
\section{Tree-Level Wilson Coefficients}
\label{sec:step12}

The KK-gluon is integrated out at the matching scale $\mu = \Mkk$, producing
four $\Delta F = 2$ four-quark operators in the BMU basis (Buras, Misiak,
Urban, hep-ph/0005183).  For a meson system involving quark flavors $i$ and
$j$:
\begin{align}
  \text{prefactor} &= \frac{1}{\Mkk^2}\,, \\[6pt]
  C_1^{\mathrm{VLL}} &= \bigl(g_L[i,j]\bigr)^2 \cdot
    \frac{\text{prefactor}}{6}\,, \\
  C_1^{\mathrm{VRR}} &= \bigl(g_R[i,j]\bigr)^2 \cdot
    \frac{\text{prefactor}}{6}\,, \\
  C_4^{\mathrm{LR}} &= -\bigl(g_L[i,j]\cdot g_R[i,j]\bigr) \cdot
    \text{prefactor}\,, \\
  C_5^{\mathrm{LR}} &= \bigl(g_L[i,j]\cdot g_R[i,j]\bigr) \cdot
    \frac{\text{prefactor}}{3}\,,
\end{align}
where $g_L[i,j]$ and $g_R[i,j]$ are the (complex) off-diagonal entries of the
coupling matrices from \hyperref[sec:step11]{Step~11}.

The Wilson coefficients are complex. Since $g_L[i,j]$ and $g_R[i,j]$ are
off-diagonal entries of Hermitian matrices, the products $g_L^2$, $g_R^2$, and
$g_L\cdot g_R$ are generally complex numbers. The imaginary parts carry
CP-violating information and are essential for the $\epsilon_K$ constraint.

The numerical coefficients come from the
color algebra of tree-level single-KK-gluon exchange in the $t$-channel.  Each
vertex carries a color generator $T^a$, giving:
\begin{equation}
  (T^a)_{\alpha\beta}\,(T^a)_{\gamma\delta} = \tfrac{1}{2}\,
    \delta_{\alpha\delta}\,\delta_{\gamma\beta}
    - \tfrac{1}{2N_c}\,\delta_{\alpha\beta}\,\delta_{\gamma\delta}
\end{equation}
After Fierz rearrangement into the BMU basis, the VLL/VRR operators pick up
$1/(2N_c) = 1/6$ for $N_c = 3$. The LR operators arise from
Fierz-rearranging the $(V{-}A)\times(V{+}A)$ structure into
scalar/pseudoscalar operator pairs; the specific coefficients $-1$ and $+1/3$
for $C_4^{\mathrm{LR}}$ and $C_5^{\mathrm{LR}}$ follow from the Fierz
identity for the LR color contractions (see e.g., Buras, Weak Hamiltonian
review, Eq.~6.6).

Tree-level single vector boson exchange
produces only vector current structures at the vertices ($\gamma^\mu P_L$ or
$P_R$). This gives VLL, VRR, and LR operators but not scalar LL/RR operators
(SLL, SRR), which would require scalar current insertions.  The full BMU basis
has 8 operators, but SLL/SRR are absent at tree level.  They could appear at
one loop but are not included here.

The flavor indices for each meson system are:

\begin{table}[h]
\centering
\begin{tabular}{@{}llll@{}}
\toprule
System & Sector & $(i,j)$ & Coupling matrices used \\
\midrule
$\epsilon_K$ & down & $(0,1) = d,s$ &
  $g_L^{\mathrm{down}},\; g_R^{\mathrm{down}}$ \\
$\Delta M_K$ & down & $(0,1) = d,s$ &
  $g_L^{\mathrm{down}},\; g_R^{\mathrm{down}}$ \\
$\Delta M_{B_d}$ & down & $(0,2) = d,b$ &
  $g_L^{\mathrm{down}},\; g_R^{\mathrm{down}}$ \\
$\Delta M_{B_s}$ & down & $(1,2) = s,b$ &
  $g_L^{\mathrm{down}},\; g_R^{\mathrm{down}}$ \\
$\Delta M_{D^0}$ & up & $(0,1) = u,c$ &
  $g_L^{\mathrm{up}},\; g_R^{\mathrm{up}}$ \\
\bottomrule
\end{tabular}
\end{table}

$\epsilon_K$ and $\Delta M_K$ use the same Wilson coefficients (same quark
pair); they differ only in which observable is computed from $M_{12}$
(\hyperref[sec:step15]{Step~15}).

Assumptions: we use tree-level matching only (no one-loop corrections at the
matching scale, no threshold corrections, no scheme conversion beyond tree
level), include only the first KK-gluon mode (higher modes are not summed,
typically a $\sim 10$--$20\%$ correction), and take $\Mkk = \LamIR$ as both
the propagator mass and the matching scale. The physical first KK-gluon mass
$m_g^{(1)} = 2.45\times\LamIR$ is corrected for at plotting time
(\hyperref[sec:step16b]{Step~16b}).

%% ====================================================================
\section{QCD Running}
\label{sec:step13}

The Wilson coefficients are evaluated at the matching scale $\Mkk$ (multi-TeV).
The hadronic matrix elements (\hyperref[sec:step14]{Step~14}) are evaluated at
$\mu_{\mathrm{had}} = 2$~GeV.  The Wilson coefficients must be evolved from
$\Mkk$ down to $\mu_{\mathrm{had}}$ using the QCD renormalization group.

\subsection{13a: VLL and VRR operators (multiplicative running)}
\label{sec:step13a}

The VLL and VRR operators do not mix with any other operators under QCD.  Their
anomalous dimension is
$\gamma_{\mathrm{VLL}} = 6(N_c-1)/N_c = 4$ for $N_c = 3$.

For a single segment with $n_f$ active flavors between scales $\mu_{\mathrm{high}}$
and $\mu_{\mathrm{low}}$:
\begin{equation}
  C_{\mathrm{VLL}}(\mu_{\mathrm{low}}) =
    C_{\mathrm{VLL}}(\mu_{\mathrm{high}}) \cdot
    \left[\frac{\alpha_s(\mu_{\mathrm{low}})}
               {\alpha_s(\mu_{\mathrm{high}})}\right]^{\gamma_{\mathrm{VLL}}/(2\beta_0)}
\end{equation}
where $\beta_0 = (33 - 2n_f)/3$ is the one-loop beta function coefficient.

Since $\gamma_{\mathrm{VLL}} > 0$ and
$\alpha_s(\mu_{\mathrm{low}}) > \alpha_s(\mu_{\mathrm{high}})$, the ratio
$> 1$ and the VLL coefficient is enhanced at lower scales.

\subsection{13b: LR operators ($2\times2$ matrix mixing)}
\label{sec:step13b}

$C_4^{\mathrm{LR}}$ and $C_5^{\mathrm{LR}}$ mix under QCD running.  The
anomalous dimension matrix is:
\begin{equation}
  \gamma_{\mathrm{LR}} = \begin{pmatrix} 8 & -4 \\ -16/3 & -28/3 \end{pmatrix}
\end{equation}

For each $n_f$ segment, the evolution matrix is:
\begin{equation}
  U_{\mathrm{LR}} = M \cdot \operatorname{diag}\!\left(
    \left[\frac{\alpha_s^{\mathrm{low}}}{\alpha_s^{\mathrm{high}}}\right]
    ^{d_i/(2\beta_0)}\right) \cdot M^{-1}
\end{equation}
where $d_i$ are the eigenvalues of $\gamma_{\mathrm{LR}}$ and $M$ is the
eigenvector matrix.  The dominant eigenvalue is large and positive, giving a
significant enhancement (${\sim}\,3$--$5\times$) of the LR operators at low
scales.  Combined with the chirality-enhanced kaon matrix elements
(proportional to $(m_K/(m_s + m_d))^2 \sim 25$), this QCD amplification is
why $\epsilon_K$ is such a powerful constraint in RS models.

\subsection{13c: Multi-segment evolution with flavor thresholds}
\label{sec:step13c}

The running from $\Mkk$ to $\mu_{\mathrm{had}} = 2$~GeV crosses one quark
threshold ($m_b = 4.18$~GeV; $m_c = 1.27$~GeV is below $\mu_{\mathrm{had}}$):
\begin{align*}
  \Mkk &\longrightarrow m_b = 4.18\;\text{GeV} \quad (n_f = 5) \\
  m_b  &\longrightarrow 2.0\;\text{GeV} \quad (n_f = 4)
\end{align*}

At each threshold, the Wilson coefficients are matched continuously (no finite
matching corrections at leading order).

The $\alpha_s$ running within each segment uses a one-loop formula:
\begin{equation}
  \alpha_s(\mu) = \frac{\alpha_s(\mu_{\mathrm{ref}})}
    {1 + \dfrac{\beta_0\,\alpha_s(\mu_{\mathrm{ref}})}{2\pi}\,
    \ln\!\left(\dfrac{\mu}{\mu_{\mathrm{ref}}}\right)}
\end{equation}
with $\alpha_s(M_Z = 91.19\;\text{GeV}) = 0.1179$ as the reference value, and
threshold matching at $m_b$ and $m_c$.

The cumulative evolution factors for VLL/VRR are multiplied across segments;
the $2\times2$ LR matrices are multiplied in order.

Assumptions: we use leading-log anomalous dimensions only (one-loop $\gamma$),
which is adequate for the current precision. The $\alpha_s$ running for the
operator evolution uses a one-loop formula ($\alpha_s(M_Z) = 0.1179$), while
the coupling computation (\hyperref[sec:step11]{Step~11}) uses 4-loop running
($\alpha_s(M_Z) = 0.1180$); this $O(\text{few-\%})$ inconsistency is
subdominant compared to the $\gs$ uncertainty. Threshold matching is continuous
(no finite corrections at $m_b$ or $m_c$), following the standard LO
prescription.

%% ====================================================================
\section{Hadronic Matrix Elements}
\label{sec:step14}

The RG-evolved Wilson coefficients at $\mu_{\mathrm{had}} = 2$~GeV are
contracted with hadronic matrix elements to obtain $M_{12}^{\mathrm{NP}}$, the
new-physics contribution to the meson mixing amplitude.

The general formula for a pseudoscalar meson $P$ with valence quarks $q_1$,
$q_2$ is:
\begin{equation}
  M_{12}^{\mathrm{NP}} = C_1^{\mathrm{VLL}} \cdot \mathrm{ME}_{\mathrm{VLL}}
    + C_1^{\mathrm{VRR}} \cdot \mathrm{ME}_{\mathrm{VLL}}
    + C_4^{\mathrm{LR}} \cdot \mathrm{ME}_{\mathrm{LR4}}
    + C_5^{\mathrm{LR}} \cdot \mathrm{ME}_{\mathrm{LR5}}
\end{equation}

VRR has the same matrix element as VLL by parity symmetry of QCD:
$\langle\bar{P}|\mathcal{O}_{\mathrm{VRR}}|P\rangle =
\langle\bar{P}|\mathcal{O}_{\mathrm{VLL}}|P\rangle$.

The matrix elements are:
\begin{align}
  \mathrm{ME}_{\mathrm{VLL}} &= \tfrac{2}{3}\,f_P^2\,m_P\,B_1\,, \\[4pt]
  r_{\chi} &= \left(\frac{m_P}{m_{q_1} + m_{q_2}}\right)^2\,, \\[4pt]
  \mathrm{ME}_{\mathrm{LR4}} &= \bigl(\tfrac{r_{\chi}}{6}
    + \tfrac{1}{4}\bigr)\,f_P^2\,m_P\,B_4\,, \\[4pt]
  \mathrm{ME}_{\mathrm{LR5}} &= \bigl(\tfrac{r_{\chi}}{2}
    + \tfrac{1}{12}\bigr)\,f_P^2\,m_P\,B_5\,.
\end{align}

The factor $r_{\chi}$ is the chiral enhancement factor.  For kaons:
$r_{\chi} = (0.49761/(0.0934 + 0.00467))^2 \approx 25.2$.
The LR operators dominate $\epsilon_K$ because their matrix elements are
enhanced by ${\sim}\,25$ relative to the VLL matrix element.

\subsection{Hadronic inputs by meson system}

\noindent\emph{Kaon ($K^0$), valence quarks $d$, $s$:}

\begin{table}[h]
\centering
\begin{tabular}{@{}lll@{}}
\toprule
Input & Value & Source \\
\midrule
$f_K$ & 155.7~MeV & PDG 2024 \\
$m_K$ & 497.61~MeV & PDG 2024 \\
$m_s(2\;\text{GeV})$ & 93.4~MeV & FLAG 2024 \\
$m_d(2\;\text{GeV})$ & 4.67~MeV & FLAG 2024 \\
$B_1^K$ & 0.717 & FLAG/lattice average \\
$B_4^K$ & 0.78 & ETM 2013 \\
$B_5^K$ & 0.57 & ETM 2013 \\
\bottomrule
\end{tabular}
\end{table}

\noindent\emph{$B_d$, valence quarks $d$, $b$:}

\begin{table}[h]
\centering
\begin{tabular}{@{}lll@{}}
\toprule
Input & Value & Source \\
\midrule
$f_{B_d}$ & 190.0~MeV & FLAG 2024 \\
$m_{B_d}$ & 5279.72~MeV & PDG 2024 \\
$m_b(m_b)$ & 4.18~GeV & PDG ($\overline{\mathrm{MS}}$) \\
$m_d(2\;\text{GeV})$ & 4.67~MeV & FLAG 2024 \\
$B_1^{B_d}$ & 0.87 & FLAG 2024 \\
$B_4^{B_d}$ & 1.02 & FLAG 2024 \\
$B_5^{B_d}$ & 0.96 & FLAG 2024 \\
\bottomrule
\end{tabular}
\end{table}

\noindent\emph{$B_s$, valence quarks $s$, $b$:}

\begin{table}[h]
\centering
\begin{tabular}{@{}lll@{}}
\toprule
Input & Value & Source \\
\midrule
$f_{B_s}$ & 230.3~MeV & FLAG 2024 \\
$m_{B_s}$ & 5366.92~MeV & PDG 2024 \\
$m_s(2\;\text{GeV})$ & 93.4~MeV & FLAG 2024 \\
$B_1^{B_s}$ & 0.87 & FLAG 2024 \\
$B_4^{B_s}$ & 1.02 & FLAG 2024 \\
$B_5^{B_s}$ & 0.96 & FLAG 2024 \\
\bottomrule
\end{tabular}
\end{table}

\noindent\emph{$D^0$, valence quarks $u$, $c$:}

\begin{table}[h]
\centering
\begin{tabular}{@{}lll@{}}
\toprule
Input & Value & Source \\
\midrule
$f_D$ & 212.0~MeV & FLAG 2024 \\
$m_{D^0}$ & 1864.84~MeV & PDG 2024 \\
$m_c(m_c)$ & 1.27~GeV & PDG ($\overline{\mathrm{MS}}$) \\
$m_u(2\;\text{GeV})$ & 2.16~MeV & FLAG 2024 \\
$B_1^D$ & 0.75 & Lattice (less precise) \\
$B_4^D$ & 1.0 & Estimated \\
$B_5^D$ & 1.0 & Estimated \\
\bottomrule
\end{tabular}
\end{table}

Assumptions: bag parameters $B_4^D$ and $B_5^D$ for $D^0$ mixing are set
to~$1.0$ because no reliable lattice determination exists, introducing $O(1)$
uncertainty in the $D^0$ constraint. The quark masses entering the chiral
enhancement factor $r_{\chi}$ are $\overline{\mathrm{MS}}$ masses at their
respective self-scales, not all at $\mu = 2$~GeV. For $B$~mesons,
$m_b(m_b) = 4.18$~GeV is used rather than $m_b(2\;\text{GeV}) \sim 5.4$~GeV;
for $D^0$, $m_c(m_c) = 1.27$~GeV. Using $m_b(2\;\text{GeV})$ instead would
change $r_{\chi}$ for $B_d$ from ${\sim}\,1.6$ to ${\sim}\,1.0$, an
$O(40\%)$ shift in the LR matrix element. For the $B$ and $D$ systems the LR
contribution is subdominant ($r_{\chi} \sim 1$--$2$, vs ${\sim}\,25$ for
kaons), so this inconsistency has limited impact on the final bounds, but it
is a known systematic.

The matrix element formulae use the vacuum insertion approximation corrected by
the bag parameters $B_i$, which encode deviations from naive factorization
$\langle P|\mathcal{O}|P\rangle =
\langle P|\bar{q}q|0\rangle\langle 0|\bar{q}q|P\rangle$.
Lattice QCD determines the $B_i$.

$M_{12}^{\mathrm{NP}}$ as computed here
is the full off-diagonal Hamiltonian matrix element
$\langle\bar{P}|\mathcal{H}_{\mathrm{eff}}|P\rangle$, with dimensions of GeV.
It does not include a $1/(2m_P)$ normalization that some references fold into
the definition of~$M_{12}$.  Consistently, the mass difference is
$\Delta m = 2|M_{12}|$, so the constraint
$|M_{12}^{\mathrm{NP}}| \le \Delta m/2$ is equivalent to
$|M_{12}^{\mathrm{NP}}| \le |M_{12}^{\mathrm{exp}}|$.

%% ====================================================================
\section{Observable Evaluation and Acceptance}
\label{sec:step15}

\subsection{15a: $\epsilon_K$ (CP violation in kaon mixing)}
\label{sec:step15a}

The NP contribution to $\epsilon_K$ is:
\begin{equation}
  \epsilon_K^{\mathrm{NP}} = \frac{|\kappa_\epsilon|}
    {\sqrt{2}\,\Delta m_K} \cdot |\mathrm{Im}(M_{12}^{\mathrm{NP}})|
\end{equation}
with $\kappa_\epsilon = 0.94$ (accounts for long-distance and higher-order
corrections, Buras et al.) and
$\Delta m_K = 3.484\times10^{-15}$~GeV ($K_L$--$K_S$ mass difference, PDG).

The NP budget is the gap between the experimental value and the SM prediction:
\begin{equation}
  \text{budget} = |\epsilon_K^{\mathrm{exp}} - \epsilon_K^{\mathrm{SM}}|
    = |2.228\times10^{-3} - 1.81\times10^{-3}| = 4.18\times10^{-4}
\end{equation}

Ratio to bound: $\epsilon_K^{\mathrm{NP}} / \text{budget}$.
Passes if ratio $\le 1.0$.

We take the SM prediction $\epsilon_K^{\mathrm{SM}} = 1.81\times10^{-3}$ from
the CKMfitter central value. The experimental value is $2.228\times10^{-3}$
(PDG), leaving a difference of $4.18\times10^{-4}$ as the allowed room for NP.
Some analyses use a tighter SM prediction, which would shrink the budget.

The formula uses $|\mathrm{Im}(M_{12}^{\mathrm{NP}})|$, discarding the sign of
the NP CP-violating phase. We do not account for constructive vs destructive
interference between NP and SM: if the NP phase partially cancels the SM
contribution the actual constraint would be weaker, and if it enhances it,
stronger. Using the absolute value is the standard simplification for
model-independent bounds.

\subsection{15b: $\Delta M_K$ (kaon mass splitting)}
\label{sec:step15b}

The NP contribution to the kaon mass difference is:
\begin{equation}
  \frac{|M_{12}^{\mathrm{NP}}|_K}{\Delta m_K / 2} = \text{ratio}
\end{equation}

Budget $= \Delta m_K/2 = 1.742\times10^{-15}$~GeV.

The constraint requires $|M_{12}^{\mathrm{NP}}|$ not to exceed
$|M_{12}^{\mathrm{exp}}| = \Delta m_K/2$.  This is conservative because the SM
contribution to $\Delta M_K$ is dominated by poorly-known long-distance effects,
so the entire experimental value is treated as a budget for NP (rather than
requiring $|M_{12}^{\mathrm{NP}}| < |M_{12}^{\mathrm{exp}} -
M_{12}^{\mathrm{SM}}|$, which would need a reliable SM prediction).
Passes if ratio $\le 1.0$.

\subsection{15c: $B_d$, $B_s$, $D^0$ mixing (mass differences)}
\label{sec:step15c}

For each system, $M_{12}^{\mathrm{NP}}$ is computed using the system-specific
hadronic parameters.  The budget definitions are:

\begin{table}[h]
\centering
\begin{tabular}{@{}lll@{}}
\toprule
System & Budget formula & Numerical value \\
\midrule
$B_d$ & $\max(\Delta m^{\mathrm{exp}}/2,\;
  |\Delta m^{\mathrm{exp}} - \Delta m^{\mathrm{SM}}|/2)$
  & $\max(1.667\times10^{-13},\;1.33\times10^{-14})$ \\
  & & $= \mathbf{1.667\times10^{-13}}$~GeV \\[4pt]
$B_s$ & $\max(\Delta m^{\mathrm{exp}}/2,\;
  |\Delta m^{\mathrm{exp}} - \Delta m^{\mathrm{SM}}|/2)$
  & $\max(5.844\times10^{-12},\;6.0\times10^{-15})$ \\
  & & $= \mathbf{5.844\times10^{-12}}$~GeV \\[4pt]
$D^0$ & $\Delta m^{\mathrm{exp}}/2$ (SM long-distance dominated)
  & $\mathbf{3.125\times10^{-15}}$~GeV \\
\bottomrule
\end{tabular}
\end{table}

The $\max(\ldots)$ formula is conservative: it uses whichever is larger,
half the experimental value or half the experiment--SM difference.  For $B_d$
and $B_s$, the experimental value dominates (the SM predictions are close to
experiment).

Ratio to bound: $|M_{12}^{\mathrm{NP}}| / \text{budget}$.
Passes if ratio $\le 1.0$.

\subsection{15d: Experimental inputs for budgets}
\label{sec:step15d}

\begin{table}[h]
\centering
\begin{tabular}{@{}lll@{}}
\toprule
System & $\Delta m^{\mathrm{exp}}$ & $\Delta m^{\mathrm{SM}}$ \\
\midrule
$K$ & $3.484\times10^{-15}$~GeV & (enters only via $\epsilon_K$) \\
$B_d$ & $3.334\times10^{-13}$~GeV & $3.6\times10^{-13}$~GeV (CKMfitter) \\
$B_s$ & $1.1688\times10^{-11}$~GeV & $1.17\times10^{-11}$~GeV (CKMfitter) \\
$D^0$ & $6.25\times10^{-15}$~GeV & Long-distance dominated (HFLAV) \\
\bottomrule
\end{tabular}
\end{table}

\subsection{15e: Overall acceptance}
\label{sec:step15e}

A scan point is accepted if the ratio to bound is $\le 1.0$ for all five
systems simultaneously: $\epsilon_K$, $\Delta M_K$, $\Delta M_{B_d}$,
$\Delta M_{B_s}$, $\Delta M_{D^0}$. These are deliberately conservative
acceptance rules (generous budgets, NP-only with no SM interference), so the
resulting bounds should be read as a weak ``best-case NP-only envelope,'' not
as a sharp physical exclusion in a global-fit sense.

Additionally, the fit must have converged with $\text{fit\_score} \le 0.1$
(meaning the quark masses and CKM are reproduced well enough that the Wilson
coefficients are meaningful).

$\Delta M_K$ was added as a 5th constraint system in the scan pipeline used for
the publication figures; earlier versions used only the other four.

%% ====================================================================
\section{Post-Processing for Figures}
\label{sec:step16}

The raw scan stores results in JSONL format (one JSON object per scan point).
Before plotting, several transformations are applied.

\subsection{16a: Quality filter}
\label{sec:step16a}

Points where the fit did not converge or where $\text{fit\_score} > 0.1$ are
discarded.  A poor fit means the quark masses and CKM were not reproduced, so
the Wilson coefficients are unreliable.

\subsection{16b: Mass-axis convention ($\LamIR \to m_g^{(1)}$)}
\label{sec:step16b}

The scan uses $\Mkk = \LamIR$ ($\xi_{\mathrm{KK}} = 1$) as the propagator mass
in all Wilson coefficient formulas (\hyperref[sec:step12]{Step~12}).  The
physical first KK-gluon mass is larger:
\begin{equation}
  m_g^{(1)} = \xi_{\mathrm{kk}} \cdot \LamIR
\end{equation}
where $\xi_{\mathrm{kk}} = 2.448687$ is the first root of the gauge-sector KK
mass eigenvalue equation with Neumann--Neumann boundary conditions on a warped
interval.

To correctly relabel the axis, we must account for the fact that the scan used
$\LamIR$ (not $m_g^{(1)}$) in the $1/\Mkk^2$ propagator.  Since the Wilson
coefficients scale as $1/\Mkk^2$, a point at $\LamIR$ with ratio~$R$ should
have ratio $R/\xi_{\mathrm{kk}}^2$ when interpreted as being at
$m_g^{(1)} = \xi_{\mathrm{kk}}\cdot\LamIR$:
\begin{align}
  \text{axis\_scale} &= \xi_{\mathrm{kk}} / \xi_{\mathrm{scan}}
    = 2.448687 / 1.0 = 2.448687\,, \\
  \Mkk^{\mathrm{plot}} &= \LamIR \cdot \text{axis\_scale}\,, \\
  \text{ratio}^{\mathrm{plot}} &= \text{ratio}^{\mathrm{stored}} /
    \text{axis\_scale}^2\,.
\end{align}

The net effect: a point that was marginal (ratio $= 1$) in the raw scan becomes
well within bounds (ratio $= 1/\xi_{\mathrm{kk}}^2 \sim 1/6$) after the
correction, because the physical propagator mass is $2.45\times$ larger than
what the scan used.  The exclusion boundary shifts to a lower $\LamIR$ (where
the raw ratio was ${\sim}\,\xi_{\mathrm{kk}}^2 \sim 6$).  In terms of
$m_g^{(1)}$, this lower $\LamIR$ multiplied by $\xi_{\mathrm{kk}}$ partially
compensates, so the boundary $m_g^{(1)}$ is similar to (but not exactly equal
to) the raw-scan $\LamIR$ boundary.

\subsection{16c: Grid aggregation (common best point)}
\label{sec:step16c}

At each unique $(r,\,\Mkk)$ grid point, there are up to 20 different
\texttt{overall\_scale} values (from the scan grid).  The plotter selects the
single \texttt{overall\_scale} that minimizes:
\begin{equation}
  \max\!\bigl(\text{ratio}_{\epsilon_K},\;\text{ratio}_K,\;
    \text{ratio}_{B_d},\;\text{ratio}_{B_s},\;\text{ratio}_{D^0}\bigr)
\end{equation}
across all five systems simultaneously.  All five system ratios are then read
from that same point.

All five contour lines therefore come from one consistent MFV point at each
grid location.  Sweeping \texttt{overall\_scale} at fixed $(r,\,\LamIR)$
explores different tradeoffs between reproducing the diagonal quark masses and
minimizing off-diagonal FCNCs; picking the best one finds the most favorable
tradeoff.

%% ====================================================================
\section{Figure 1: Exclusion Boundaries}
\label{sec:step17}

Figure~1 shows contour lines in the $(r,\,m_g^{(1)})$ plane where each meson
system's $\text{ratio\_to\_bound} = 1$, with green/red shading for the combined
allowed/excluded regions.

The aggregated data from \hyperref[sec:step16c]{Step~16c} is interpolated onto a
regular grid in $(\log_{10} r,\;\log_{10} \Mkk)$ space.  For each system, a
contour is drawn at $\text{ratio} = 1$.  The region above all contours is shaded
green (allowed); below is red (excluded).

At a given $r$, any $m_g^{(1)}$ above all contour lines is allowed by all five
meson constraints.  The contour that sits highest at each $r$ is the binding
constraint.  Typically $\epsilon_K$ dominates for $r \sim 0.1$--$0.5$, and
$\Delta M_K$ takes over at large~$r$.

This is a best-case envelope.  At each grid point, the fitter found the single
most favorable Yukawa structure, and the plotter picked the most favorable
\texttt{overall\_scale}.  The contour shows the lowest $\Mkk$ at which any
fitted point survives; it is an existence bound, not a statement about typical
NMFV behavior.

\textcolor{red}{\textbf{To do:} These figures answer ``can NMFV ever allow
$\Mkk$ this low?''\ via optimization.  The analysis Jackie's note asks for is
different: generate an ensemble of random $O(1)$ NMFV Yukawas in the explicit
$(\alpha_Q,\,\beta_Q,\,\gamma_Q)$ parametrization, keep only points that
reproduce observed quark masses, and histogram the $\Mkk$ values at which
$\mathrm{Im}(C_4^K)$ meets its bound.  That would answer ``what does NMFV
\emph{typically} require?''\ Both analyses are needed.}

%% ====================================================================
\section{Figure 2: 2007 vs Modern $\Mkk$ Bound}
\label{sec:step18}

Figure~2 shows, for each $r$ value, the minimum $m_g^{(1)}$ among all accepted
points, comparing modern (2024+) constraints to rescaled 2007-era constraints.

The construction proceeds as follows:
\begin{enumerate}
  \item For each unique $r$ value in the scan grid, collect all
    $(\Mkk,\;\text{ratios})$ pairs across all $\LamIR$ and
    \texttt{overall\_scale} values.
  \item For the modern bound: check each pair's
    $\max(\text{all 5 ratios}) \le 1$.  Record the smallest $\Mkk$ among
    passing points.  If no point passes at any $\Mkk$, mark this $r$ as fully
    excluded.
  \item For the 2007 bound: multiply each system's ratio by the
    corresponding \texttt{BOUND\_RATIO} factor (which makes the constraint
    looser, since modern bounds are tighter).  Then apply the same acceptance
    check and find the minimum passing $\Mkk$.
  \item Plot both curves as a function of~$r$.
  \item Shade the region between the two curves to visualize the improvement
    since 2007.
\end{enumerate}

Figure~2 checks $\max(\text{all 5 ratios}) \le 1$ per scan row, so it does not
mix different \texttt{overall\_scale} values across systems. Unlike Figure~1, it
does not first collapse each $(r,\,\Mkk)$ to a single best-scale
representative; it considers every stored row individually.

The 2007 rescaling factors are:

These factors represent the ratio of the modern experimental bound to the
2007-era bound.  Multiplying the modern ratio by this factor gives the
2007-equivalent ratio:

\begin{table}[h]
\centering
\begin{tabular}{@{}llll@{}}
\toprule
System & Factor & Origin & Interpretation \\
\midrule
$\epsilon_K$ & 0.70 & Tighter lattice + CKM inputs & Modern 30\% tighter \\
$\Delta M_K$ & 0.70 & Same improvement as $\epsilon_K$ & Modern 30\% tighter \\
$\Delta M_{B_d}$ & 0.67 & Better lattice + exp & Modern 33\% tighter \\
$\Delta M_{B_s}$ & 0.37 & Huge improvement in $B_s$ measurements &
  Modern 63\% tighter \\
$\Delta M_{D^0}$ & 0.106 & $D^0$ mixing barely measured in 2007 &
  Modern ${\sim}\,90\%$ tighter \\
\bottomrule
\end{tabular}
\end{table}

For example, if a point has modern $B_s$ ratio $= 0.9$ (barely passing), its
2007 equivalent is $0.9\times0.37 = 0.33$ (easily passing), so the modern
constraint is much more stringent for $B_s$ mixing.

The purple shaded band between the two curves shows how much constraints have
tightened since 2007.  At small $r$ (where $\epsilon_K$ dominates), the
improvement is modest (${\sim}\,30\%$).  At large~$r$ (where $B_s$ and $D^0$
dominate), the improvement is dramatic.

%% ====================================================================
\section{Summary of Key Assumptions}
\label{sec:summary}

\begin{enumerate}
  \item Slice of $\mathrm{AdS}_5$ with
    $k = \Mpl = 1.2209\times10^{19}$~GeV, Higgs on the IR brane, no
    backreaction.

  \item SVD Yukawa parametrization with right rotations fixed to identity,
    forcing $C_u$ and $C_d$ diagonal. Only $C_Q$ has off-diagonal structure.

  \item $C_Q = r\,Y_u Y_u^\dagger + Y_d Y_d^\dagger$ with no explicit
    $\alpha_Q$ identity term. A sigmoid maps eigenvalues to
    $c \in [0.30,\,0.72]$.

  \item 14-parameter TRF optimization with max 120 evaluations per point, from
    a single deterministic seed. Finds the best Yukawa structure, not a
    representative sample.

  \item $\gs = 3.0$ (one-loop enhanced, Gedalia et al.\ 2009), entering as
    ${\gs}^2$ in all Wilson coefficients.

  \item Tree-level matching at $\Mkk = \LamIR$, first KK-gluon mode only.

  \item LO anomalous dimensions, one-loop $\alpha_s$ running, threshold
    matching at $m_b$ and $m_c$, hadronic scale $\mu_{\mathrm{had}} = 2$~GeV.

  \item FLAG 2024 / PDG 2024 hadronic inputs. $D^0$ bag parameters $B_4$,
    $B_5$ estimated at~$1.0$. Kaon $B_4$, $B_5$ from ETM 2013.

  \item $\epsilon_K$ budget: exp--SM gap $(4.18\times10^{-4})$. $\Delta M_K$:
    half the experimental value. $B_d$, $B_s$:
    $\max(\text{exp}/2,\;|\text{exp}-\text{SM}|/2)$. $D^0$: exp$/2$.

  \item Publication figures use
    $m_g^{(1)} = 2.4487\,\LamIR$ with a $1/\xi_{\mathrm{kk}}^2$ ratio
    correction.

  \item Both figures show the most favorable point at each $(r,\,\Mkk)$:
    an existence bound, not a characterization of typical NMFV behavior.

  \item Fit targets
    ($m_u = 1.3$~MeV, $m_c = 620$~MeV, $m_t = 172$~GeV, $m_d = 2.8$~MeV,
    $m_s = 57$~MeV, $m_b = 2.86$~GeV) are repo-owned values chosen for scan
    stability, not PDG $\overline{\mathrm{MS}}$ masses at a canonical scale.
\end{enumerate}

\end{document}
